\begin{frame}{Target widths without AMVF: a simulation study}
\context{Truth supplies the track groups; a separate fit estimates the vertex position.}
\small
\textbf{Group the reconstructed tracks by their truth vertex, with $\ntrk\geq2$.}\\[.3em]
Fit reconstructed measurements. Use the truth position only to measure the residual afterwards.
\par\vspace{.2em}
\textbf{Start with a simple one-dimensional fit.}\\[.3em]
At a fixed transverse beam position, a weighted least-squares estimate is
\[
\hat z=\frac{\sum_i z_{0,i}/\sigma_{z_0,i}^2}{\sum_i1/\sigma_{z_0,i}^2}.
\]
All $z_0$ values must refer to the same beam line. Check this approximation and the residual bias. A full 3D fit needs a standalone vertex fitter and full track covariances.
\par\vspace{.2em}
\textbf{Measure the width of $\hat z-z_{\rm truth}$ versus track multiplicity.}\\[.3em]
Fit a positive curve to set the Gaussian training-target width. Validate it on separate events.
\takeaway{A fit is still needed. Truth grouping replaces AMVF finding and track assignment.}
\sourcecaption{Proposal only. Truth grouping omits association mistakes. Target width is not the final position error.}
\end{frame}
